We now prove Corollary \ref{cor:universality} on universality over the
disorder $\X$, using the results of \cite{wang2024universality}.

\begin{proof}[Proof of Corollary \ref{cor:universality}]
Let $\tilde \X \in \R^{n \times n}$ satisfy Assumption \ref{assump:X-universal}.
Let $\X \in \R^{n \times n}$ be an orthogonally-invariant matrix
satisfying Assumption \ref{assump:X} with the same bound $\Lambda_{\max}$ and
asymptotic eigenvalue distribution $\Lambda$ as $\tilde\X$.
Let $\S \in \R^{n \times n}$ be
a random diagonal sign matrix as in Corollary \ref{cor:universality}(d),
and let $\bPi \in \R^{n \times n}$ be a uniformly random permutation matrix,
where $(\S,\bPi)$ are independent of each other and of $\btheta^0,\{\b^t\}_{t
\geq 0}$. Then \cite[Proposition 2.7, Theorem 2.8]{wang2024universality}
implies that the AMP algorithm \eqref{eq:discreteAMP} applied with $\tilde\X$ has the
same dynamical mean-field limit \eqref{eq:AMPconvergence} (in the sense of almost-sure
Wasserstein-2 convergence of the coordinatewise empirical distribution of its
iterates) as applied with the matrix
\[\W=\bPi\S\tilde\X\S\bPi^\top.\]
As \eqref{eq:AMPconvergence} and the bound $\|\X\|_\op \leq \Lambda_{\max}$ are the only
properties of $\X$ used in all preceding proofs, this implies that
Theorems \ref{thm:dmft-approx}, \ref{thm:replica} and Corollaries
\ref{cor:hightemp}, \ref{cor:ising} hold equally with $\W$ in place of
$\X$.

The dynamics \eqref{eq:dynamics} driven by $\W$ have the form
\[\d\btheta^t=[f(\btheta^t)+\bPi\S\tilde\X\S\bPi^\top\btheta^t]\d t
+\sqrt{2\gamma}\,\d\b^t.\]
Define $\g^t=\W\btheta^t-\int_0^t r_g(t-s)\btheta^s \d s$ as in Theorem
\ref{thm:dmft-approx}, and let
\[\tilde\btheta^t=\bPi^\top \btheta^t,
\qquad \tilde\g^t=\bPi^\top \g^t,
\qquad \tilde\b^t=\bPi^\top \b^t.\]
The conclusion of Theorem \ref{thm:dmft-approx} 
for $\{(\btheta^t,\g^t,\b^t)\}_{t \geq 0}$ implies the same convergence for
$\{(\tilde\btheta^t,\tilde\g^t,\tilde\b^t)\}_{t \geq 0}$, since this has the
same empirical distribution of sample paths.
Multiplying the above dynamics by $\bPi^\top$ and noting that
$\bPi^\top f(\btheta^t)=f(\tilde\btheta^t)$ since $f:\R \to \R$ is applied
entrywise, the dynamics of $\{(\tilde\btheta^t,\tilde\g^t,\tilde\b^t)\}_{t \geq
0}$ are given by
\begin{equation}\label{eq:permutedprocess}
\d\tilde\btheta^t=[f(\tilde\btheta^t)+\S\tilde\X\S\tilde\btheta^t]\d t
+\sqrt{2\gamma}\,\d\tilde \b^t,
\qquad \tilde \g^t=\S\tilde\X\S\tilde\btheta^t-\int_0^t
r_g(t-s)\tilde\btheta^s\d s,
\end{equation}
coinciding with the dynamics \eqref{eq:dynamics} driven by $\S\tilde\X\S$.
This shows Corollary \ref{cor:universality}(d) for the dynamics.
Since the Gibbs measures $\mu_\Gibbs$ and $\mu_{\text{Ising}}$ defined by
$\W$ are also identical to those defined by $\S\tilde\X\S$,
Corollary \ref{cor:universality}(d) holds also for the Gibbs measures.

If $f(\cdot)$ is odd and $\theta^0$ is sign-invariant in law, then 
the joint law of $\{(\theta^t,g^t,b^t)\}_{t \geq 0}$ in the dynamical
mean-field limit of Definition \ref{def:DMFT} is also invariant under
multiplication by a single sign $s \in \{\pm 1\}$. Setting
\[\tilde\btheta^t=\S\bPi^\top \btheta^t,
\qquad \tilde\g^t=\S\bPi^\top \g^t,
\qquad \tilde\b^t=\S\bPi^\top \b^t,\]
the conclusion of Theorem \ref{thm:dmft-approx} 
for $\{(\btheta^t,\g^t,\b^t)\}_{t \geq 0}$ implies that the empirical
distribution of sample paths of
$\{(\tilde\btheta^t,\tilde\g^t,\tilde\b^t)\}_{t \geq 0}$ converges to
the law of $\{(s\theta^t,sg^t,sb^t)\}_{t \geq 0}$ where $s \in \{\pm 1\}$ is
an independent and uniformly random sign. This is the same law as
that of $\{(\theta^t,g^t,b^t)\}_{t \geq 0}$, so Theorem \ref{thm:dmft-approx} 
holds also for $\{(\tilde\btheta^t,\tilde\g^t,\tilde\b^t)\}_{t \geq 0}$.
Since $\S\bPi^\top f(\btheta^t)=f(\S\bPi^\top\btheta^t)=f(\tilde\btheta^t)$ when
$f(\cdot)$ is odd, we have
\[\d\tilde\btheta^t=[f(\tilde\btheta^t)+\tilde\X\tilde\btheta^t]\d t
+\sqrt{2\gamma}\,\d\tilde \b^t,
\qquad \tilde \g^t=\tilde\X\tilde\btheta^t-\int_0^t
r_g(t-s)\tilde\btheta^s\d s\]
coinciding with the dynamics \eqref{eq:dynamics} driven by $\tilde\X$.
This shows
Corollary \ref{cor:universality}(a). If $U(\cdot)$ is even and $h=0$ in the
Ising model, then the Gibbs measures $\mu_\Gibbs$ and $\mu_\text{Ising}$ defined
by $\W$ are also identical to those defined by $\tilde \X$, implying
Corollary \ref{cor:universality}(b--c).
\end{proof}
